\subsection{Spielstruktur}\label{subsec:spielstruktur}


"`Breaking Worse"' ist ein Spiel, in dem die Spielenden in die Welt der Drogen eintauchen und sich mit Aktivitäten beschäftigen, die von der Herstellung und dem Kochen bis hin zur Verteilung und dem Verkauf reichen. Die Spielfiguren stehen vor einer ernsten gesundheitlichen Situation, die sie dazu zwingt, ein teures Heilmittel zu kaufen, um die tödliche Krankheit zu besiegen, bevor die Zeit abläuft.

\subsubsection{Spielphasen}

\paragraph{Early-Game} Die Spielfiguren beginnen in einer Küche,
 in der ihnen verschiedene Werkzeuge und Rezepte zur Verfügung stehen.
 Um mit dem Kochen zu beginnen, müssen sie jedoch nach draußen gehen,
 um auf der Karte wichtige Zutaten zu sammeln. Sobald sie diese Zutaten eingesammelt haben, 
 kehren sie in die Küche zurück, um mit dem Herstellen zu beginnen.

\paragraph{Mid-Game} Nach dem Fertigstellen der Kreationen können die Spieler:innen erneut die Karte erkunden, 
um ihre Produkte zu verkaufen. Dies geschieht durch die Interaktion mit NPCs, 
die jeweils bestimmte Bedürfnisse und Mengen haben. 
Dabei müssen die Spieler:innen aufmerksam bleiben, denn die Stadt wird streng von der Polizei überwacht. 
Werden sie erwischt, drohen ihnen Zeit- und Geldstrafen.

\paragraph{Late-Game} Während die Spielfiguren weiter herstellen und verkaufen,
steigt ihr Verdachtslevel, was erhöhte Vorsicht erfordert und das Risiko höherer Zeitstrafen
durch Gefängnisaufenthalte mit sich bringt. Mit ausreichend Geld können sie Werkzeuge aufrüsten
und neue Rezepte freischalten, um stärkere Substanzen herzustellen. Das ultimative Ziel des Spiels ist es,
den finalen Gegenstand zu erwerben: das Heilmittel gegen Krebs,
um die Krankheit endgültig zu besiegen.

\subsubsection{Dynamik des Spiels}

Einige Gebiete der Karte werden vom rivalisierenden Kartell besetzt, das nicht zögert, 
auf die Spielfiguren zu schießen, wenn diese ihren Weg kreuzen. 
Jeder Schaden, den die Spieler:innen von ihnen erleiden, 
verringert die gemeinsame Gesundheitsanzeige, 
die entweder mit der Zeit regeneriert oder gegen eine Gebühr im Krankenhaus wiederhergestellt werden kann.
 Ein tödlicher Angriff durch das Kartell lässt die Spielfiguren im Krankenhaus mit geringer Gesundheit und einer erheblichen Geldstrafe wieder erscheinen.

Die Polizeialarmbereitschaft wird direkt durch das Aktivitätsniveau der Spielfiguren beeinflusst.
Die Spielenden haben die Wahl, entweder mit Gewalt gegen die Polizei vorzugehen,
was deren Wachsamkeit und die Schwierigkeit erhöht,
oder sich vorsichtig zu verhalten,
indem sie zum Beispiel die Spielfiguren in Büschen verstecken wenn diese verfolgt werden,
um Konfrontationen zu vermeiden.

\subsubsection{Mögliche Spielmodi und Strategie}

Das Spiel bietet 3 Schwierigkeitsgrade, je nach dem gibt es weniger Zeit, schneller steigendes Fahndungslevel niedrigere Verkaufspreise und längere Gefängnisstrafe.

Zusammenfassend erlaubt es "`Breaking Worse"' den Spielenden, ihre Herangehensweise und Strategie 
flexibel zu gestalten und durch taktisches Verhalten langfristig zu profitieren.

